\renewcommand{\abstractname}{Agradecimientos}
\begin{abstract}
Incluir aquí los agradecimientos que se deseen, dedicatoria y cualquier otro texto de carácter personal.

\lipsum[1-2]
\end{abstract}

\subsection{Imágenes térmicas en el espectro infrarrojo}

    Uno de los principales retos a los que se enfrentan los sistemas de seguimiento de UAVs es la condición del entorno. Un buen sistema SOT debe poder detectar y seguir objetivos en cualquier momento del día y bajo cualquier condición adversa (temperatura, humedad, temporal...), lo cual es especialmente complicado en las horas nocturnas o de poca luz natural, y en entornos con niebla, lluvia, rachas fuertes de viento y demás fenómenos ambientales. Estos factores no solo reducen a menudo la calidad de la imagen, sino que además introducen ruido y variaciones en la firma del objetivo. Es por esto que la tecnología de imágenes infrarrojas comenzó a ser muy relevante en este campo.

    La tecnología de procesado de ímagenes térmicas en el espectro infrarrojo (TIR) consiste en la reconstrucción de imágenes a través de la radiación infrarroja emitida por los objetos de la escena. Una de las principales ventajas de este proceso es su capacidad para no interferir directamente con el objetivo, ya que la radiación electromagnética es una propiedad física inherente a cualquier cuerpo con una temperatura por encima del cero absoluto. Al limitarse a recibir dicha señal sin emitir energía externa, los detectores TIR se clasifican como sistemas de detecciñón pasivos, lo que los hace ideales para aplicaciones militares y de vigilancia. 

    Pese a todas sus ventajas, el proceso de obtención de estas imágenes requiere una serie de etapas de procesamiento debido al bajo contraste y a la escasa resolución de los detalles de la imagen. Estas etapas de procesamiento suelen llevarse a cabo a través técnicas de visión artificial (computer vision) y aprendizaje profundo (deep learning) \cite{review_infrared_imaging}.
